\chapter{Complexity — Poly Time Block Tests}
%%% generated by the blueprint pipeline from TCSlib.Complexity.ClassNP.PolyTimeBlockTests
\section{Overview}

Declarations of the Lean module \texttt{TCSlib.Complexity.ClassNP.PolyTimeBlockTests}.

\section{Declarations}

\begin{definition}[Block at a given index]
\lean{Complexity.blockAt}
\label{Complexity.blockAt}
\leanok
\uses{Complexity.pairFstD, Complexity.pairSndD}
Given natural numbers $a$ and $k$, a Boolean list $z$, and an index $i$, decode $z$ as a
pair of lists and take from the second list the block of length $a(m+1)^k$ beginning at
position $i\,a(m+1)^k$, where $m$ is the length of the first decoded list.
\end{definition}

\begin{theorem}[Length of a Boolean flag]
\lean{Complexity.BlockLoop.length_flag_le}
\label{Complexity.BlockLoop.length_flag_le}
\leanok
\difficulty{4}
If $f$ is a nonempty finite bit string and $b$ is a Boolean value, the length of the
string that is $[1]$ when $b$ is true and $f$ otherwise is at most the length of $f$.
\end{theorem}

\begin{theorem}[Success on an extended range]
\lean{Complexity.BlockLoop.any_range_succ_flag}
\label{Complexity.BlockLoop.any_range_succ_flag}
\leanok
\difficulty{4}
For any Boolean-valued function $F$ on the natural numbers and any $i$, the indicator
that either $F(i)$ holds or $F$ holds somewhere below $i$ equals the indicator that $F$
holds somewhere below $i+1$.
\end{theorem}

\begin{theorem}[Polynomial time for an iterated Boolean search]
\lean{Complexity.BlockLoop.polyTimeComputable_orLoop}
\label{Complexity.BlockLoop.polyTimeComputable_orLoop}
\leanok
\uses{Complexity.BlockLoop.any_range_succ_flag, Complexity.BlockLoop.length_flag_le, Complexity.BlockLoop.loop_output, Complexity.BlockLoop.lt_polyBudget_of_le, Complexity.BlockLoop.polyTimeComputable_countdown, Complexity.BlockLoop.step, Complexity.PolyTimeComputable, Complexity.pairFstD, Complexity.pairFstD_pairEncode, Complexity.pairSndD, Complexity.polyTimeComputable_const, Complexity.polyTimeComputable_ite, Complexity.polyTimeComputable_pairFstD, Turing.length_pairEncode, Turing.pairEncode}
\difficulty{5}
Let $t$ map bit strings to Booleans, and let $\mathrm{nxt},x,r$ map bit strings to bit
strings. Let $Q$ map a bit string and a natural number to a bit string, and let $F$ map
a bit string and a natural number to a Boolean. Fix $a',k',c\in\mathbb{N}$. Suppose the
functions $s\mapsto[t(\pi_2(\pi_2(s)))]$ and $s\mapsto\mathrm{nxt}(\pi_2(\pi_2(s)))$ are
polynomial-time computable, as are $x$ and $r$. Suppose that $|\mathrm{nxt}(q)|\le
|q|+2+c$ for every bit string $q$, that $|x(z)|\le |r(z)|$ and $r(z)=Q(z,0)$ for every
$z$, and that for every $z$ and every $i<a'(|x(z)|+1)^{k'}$, one has
$\mathrm{nxt}(Q(z,i))=Q(z,i+1)$ and $t(Q(z,i))=F(z,i)$. Then the Boolean-valued
bit-string function whose output is the one-bit list indicating whether $F(z,i)$ is true
for at least one $i<a'(|x(z)|+1)^{k'}$ is polynomial-time computable.
\end{theorem}

\begin{theorem}[Polynomial-Time Block XOR Test]
\lean{Complexity.polyTimeComputable_blockXorAnyTest}
\label{Complexity.polyTimeComputable_blockXorAnyTest}
\leanok
\uses{Complexity.BlockLoop.polyTimeComputable_orLoop, Complexity.BlockLoop.sliceDropAt_pairEncode_drop, Complexity.BlockLoop.sliceTakeAt_pairEncode_drop, Complexity.PolyTimeComputable, Complexity.PolyTimeComputable.comp, Complexity.PolyTimeComputable.pairEncode, Complexity.blockAt, Complexity.length_pair_components_le, Complexity.length_sliceDropAt_le, Complexity.pairFstD, Complexity.pairFstD_pairEncode, Complexity.pairSndD, Complexity.pairSndD_pairEncode, Complexity.polyTimeComputable_pairFstD, Complexity.polyTimeComputable_pairSndD, Complexity.polyTimeComputable_sliceDropAt, Complexity.polyTimeComputable_sliceTakeAt, Complexity.polyTimeComputable_xorD, Complexity.sliceDropAt, Complexity.sliceTakeAt, Complexity.xorD, Complexity.xorD_pairEncode, Turing.MultiTapeTM.indicator, Turing.length_pairEncode, Turing.pairEncode}
\difficulty{5}
Let $V$ be a language of finite bit strings whose Boolean membership indicator, output
as a one-bit word, is polynomial-time computable. For any $a,k,a',k'\in\mathbb{N}$, the
following function is polynomial-time computable: on input $w$, parse its first
component as a pair to obtain $x$ (using the empty word if parsing fails), and let $y$
be the second component of $w$ (again using the empty word if parsing fails). For each
of the first $a'(|x|+1)^{k'}$ blocks of length $a(|x|+1)^k$ in the parsed second
component of the inner pair, take the bitwise XOR with the corresponding initial segment
of $y$, truncating to the shorter word; the output is true exactly when at least one
resulting word, paired with $x$, belongs to $V$.
\end{theorem}

\begin{theorem}[XOR of a list with itself]
\lean{Complexity.zipWith_xor_self}
\label{Complexity.zipWith_xor_self}
\leanok
\difficulty{4}
For every finite bit string $l$, its bitwise XOR with itself is the bit string of zeros
of length $|l|$.
\end{theorem}

\begin{theorem}[XOR with a zero block]
\lean{Complexity.zipWith_xor_blockAt}
\label{Complexity.zipWith_xor_blockAt}
\leanok
\uses{Complexity.blockAt, Complexity.pairFstD}
\difficulty{6}
For natural numbers $a$, $k$, and $i$, and a bit string $z$, let $m=a(|\pi_1(z)|+1)^k$,
where $\pi_1(z)$ is the total first projection of $z$. The bitwise XOR of the length-$m$
zero word with the length-$m$ block beginning at position $im$ in the total second
projection of $z$ equals that block.
\end{theorem}

\begin{theorem}[Polynomial-Time Block Membership Test]
\lean{Complexity.polyTimeComputable_blockAnyTest}
\label{Complexity.polyTimeComputable_blockAnyTest}
\leanok
\uses{Complexity.PolyTimeComputable, Complexity.PolyTimeComputable.comp, Complexity.PolyTimeComputable.pairEncode, Complexity.blockAt, Complexity.pairFstD, Complexity.pairFstD_pairEncode, Complexity.pairSndD_pairEncode, Complexity.polyTimeComputable_blockXorAnyTest, Complexity.polyTimeComputable_id, Complexity.polyTimeComputable_pairFstD, Complexity.polyTimeComputable_polyUnary, Complexity.polyTimeComputable_xorD, Complexity.xorD_pairEncode, Complexity.zipWith_xor_blockAt, Complexity.zipWith_xor_self, Turing.MultiTapeTM.indicator, Turing.pairEncode}
\difficulty{4}
Let $V$ be a language of finite bit strings whose Boolean membership indicator, output
as a one-bit word, is polynomial-time computable. For any $a,k,a',k'\in\mathbb{N}$, the
following function is polynomial-time computable: on input $z$, let $x$ and $y$ be the
first and second components obtained by parsing $z$ as a pair (using the empty word for
a component when parsing fails); divide $y$ into consecutive blocks of length
$a(|x|+1)^k$, starting at its beginning, and form the one-bit output that is true
exactly when at least one of the first $a'(|x|+1)^{k'}$ such blocks, paired with $x$,
belongs to $V$.
\end{theorem}
